\documentclass[11pt,a4paper]{article}
\usepackage[margin=1in]{geometry}
\usepackage{amsmath,amssymb,amsthm}
\usepackage[colorlinks=true,allcolors=blue]{hyperref}

\newtheorem{theorem}{Theorem}
\newtheorem{corollary}[theorem]{Corollary}
\newtheorem{lemma}[theorem]{Lemma}
\theoremstyle{definition}
\newtheorem{definition}[theorem]{Definition}
\theoremstyle{remark}
\newtheorem{remark}[theorem]{Remark}

\newcommand{\PR}{\mathrm{PR}}
\newcommand{\ind}[1]{\mathbf{1}_{#1}}
\newcommand{\ev}[1]{\langle #1\rangle_P}
\newcommand{\F}{\mathbb{F}_2}

\title{A joint quantum measure on two copies of the CHSH scenario gives Tsirelson's bound,
and the converse of Henson's theorem fails%
\thanks{Multi-DAC. Authors in alphabetical order. Clawd and Wren are AIs. How this work was done and
checked is described in the README of the repository \texttt{Multi-DAC/releases}, where the
machine-checked proofs and the scripts for every computed figure are filed with this paper.}}
\author{Clawd \and Clayton Iggulden-Schnell \and Wren}
\date{20 November 2026}

\begin{document}
\maketitle

\begin{abstract}
In Sorkin's quantum measure theory a behaviour is admitted only if the sample space of all outcomes
carries a quantum measure that returns the observed probabilities. Henson showed that this joint
quantum measure (JQM) condition implies Consistent Exclusivity, and asked whether the converse holds.
It does not. On two copies of the CHSH scenario, with the branching measurements of Dowker, Henson and
Wallden included, we prove an identity that fixes the sum of the measures of four explicit $128$-point
sets at $2-\langle C_1C_2\rangle/4$, where $C_1,C_2$ are the CHSH expressions of the two copies.
Non-negativity of the measure then gives $|\langle C_1C_2\rangle|\le 8$ for every behaviour with a JQM.
For two copies of an isotropic box this is exactly Tsirelson's bound $2\sqrt2$, where Consistent
Exclusivity gives only $4(\sqrt{10}-1)/3\approx 2.883$. Two copies of an isotropic box have a JQM if and
only if the box is quantum, and isotropic boxes just above Tsirelson's bound obey Consistent Exclusivity
and have no JQM. The identity, the bound and the non-existence statement are machine-checked in Lean~4.
\end{abstract}

\section{Introduction}

A quantum measure \cite{Sorkin} is a non-negative set function that may violate the Kolmogorov sum rule
for two disjoint sets but obeys the corresponding rule for three. Dowker, Henson and Wallden \cite{DHW}
proposed it as a replacement for the joint probability distribution in the definition of
non-contextuality: a behaviour is in the set \emph{JQM} if a single quantum measure on the space of
all outcomes of all measurements reproduces the observed probabilities. Henson \cite{Henson} proved
that a behaviour in JQM obeys Consistent Exclusivity (CE), so that, for the CHSH scenario, two
independent copies of a behaviour with CHSH value above $2.883$ are excluded. He left open whether the
converse holds and how much stronger JQM might be. Dowker, Henson and Wallden note that JQM is not
closed under composition, that two independent PR boxes are not in JQM, and that it is open whether
their stronger, strongly positive condition follows from JQM together with closure under composition.

This paper settles the two-copy CHSH case. The result is an identity (Theorem~\ref{t:identity}): four
explicit subsets of the sample space have measures whose sum is fixed, by the decoherence conditions
alone, at $2-\langle C_1C_2\rangle/4$. Since a measure is non-negative, $\langle C_1C_2\rangle\le 8$.
For two copies of a box with CHSH value $S$ this reads $S\le 2\sqrt2$. For isotropic boxes the condition
is also sufficient (Theorem~\ref{t:iso}), while Consistent Exclusivity on two copies holds up to
$S=4(\sqrt{10}-1)/3$ (Theorem~\ref{t:lo}, which reproduces the value of Sainz et al.\ \cite{Sainz} and
proves it is exact on the whole isotropic line). In the window between, a behaviour obeys CE and has no
JQM (Corollary~\ref{c:converse}).

The result depends on one premise, which we state before anything else (Section~\ref{s:setting}): the
measurements of the two-copy scenario include the branching measurements of \cite{DHW}, in which a
party's choice of box may depend on outcomes already obtained by others, across the two wings. This is
the same list of measurements on which the benchmark value $2.883$ is obtained, and
Section~\ref{s:needs} shows that neither that value nor our identity survives without it.

\section{Setting}\label{s:setting}

\paragraph{Scenario.} There are four parties, $A_1,B_1$ (copy 1) and $A_2,B_2$ (copy 2). Each has two
boxes, labelled $0$ and $1$, opens one per run and obtains an outcome $\pm1$. The sample space is
$\Omega=\{\pm1\}^8$: a point specifies an outcome for every box of every party. We write
$\alpha_x,\beta_y,\alpha'_{x'},\beta'_{y'}$ ($x,y,x',y'\in\{0,1\}$) for the outcome of box $x$ of $A_1$,
box $y$ of $B_1$, box $x'$ of $A_2$, box $y'$ of $B_2$, as functions on $\Omega$. An \emph{event} fixes
one box and one outcome for each party; it is a subset of $\Omega$ with $16$ points, and there are
$256$ events. A \emph{behaviour} $P$ assigns a probability to every event, summing to $1$ over the $16$
events of each choice of boxes. $W$ denotes the linear span of the indicator functions of events. For
$s=(x,y,x',y')$ the correlator $\ev{\alpha_x\beta_y\alpha'_{x'}\beta'_{y'}}$ is the expectation of the
product of the four outcomes under $P$ when the boxes $s$ are opened. Two events are \emph{locally orthogonal} if some party has the same box
in both with different outcomes; this is the case exactly when they are disjoint as subsets of $\Omega$.
There are $22{,}400$ locally orthogonal pairs.

\paragraph{Premise: branching measurements.} A measurement is a procedure in which each party opens one
box, and the choice of box may depend on the outcomes already obtained by other parties, in either wing
and in either copy. This is the scenario with branching measurements of \cite[Section on branching
measurements]{DHW}. Its fine-grained outcomes are events, and by the single-branching lemma of
\cite{DHW} any two locally orthogonal events are outcomes of one such measurement. So in this scenario
two events are exclusive, in the sense of \cite{Henson}, exactly when they are locally orthogonal.

\begin{definition}[JQM]\label{d:jqm}
A \emph{joint quantum measure} for a behaviour $P$ is a real bilinear form $D$ on the space of real
functions on $\Omega$ such that
\begin{itemize}
\item[(i)] $D(\ind e,\ind e)=P(e)$ for every event $e$;
\item[(ii)] $D(\ind e,\ind{e'})=0$ for every pair of locally orthogonal events $e,e'$;
\item[(iii)] $D(\ind S,\ind S)\ge 0$ for every $S\subseteq\Omega$.
\end{itemize}
\end{definition}

The quantum measure is $\mu(S)=D(\ind S,\ind S)$. In the form given by Henson \cite{Henson}, a joint
quantum measure is a function $\mu:2^\Omega\to\mathbb{R}_{\ge0}$ that obeys the three-set sum rule and
agrees with $P$ on every coarse-grained outcome of every measurement; \cite[the lemma following the
definition of JQM]{DHW} restates this with a decoherence functional $D$ that equals $P$ on the diagonal
and vanishes on disjoint outcomes of one measurement. Under the premise this implies (i)--(iii) for the
real part of $D$, which is all we use for non-existence. For existence (Theorem~\ref{t:iso}) we construct
a strongly positive decoherence functional that satisfies the conditions of \cite{DHW} in full. Symmetry
of $D$ is never used.

\paragraph{Consistent Exclusivity.} A behaviour obeys CE if every set of pairwise exclusive events has
total probability at most $1$ \cite{Henson}. Under the premise this is local orthogonality for two
copies, $\mathrm{LO}^2$ in the notation of \cite{Sainz}.

\paragraph{Isotropic boxes.} $\PR_q=q\,\PR+(1-q)\,U$, with $\PR$ the Popescu--Rohrlich box \cite{PR}
and $U$ the uniform box, $0\le q\le1$. Its CHSH value is $4q$, and it is quantum if and only if
$q\le 1/\sqrt2$ \cite{Tsirelson}. On copy $i$ the CHSH expression is
$C_1=\sum_{x,y}(-1)^{xy}\alpha_x\beta_y$, $C_2=\sum_{x',y'}(-1)^{x'y'}\alpha'_{x'}\beta'_{y'}$, and
$\ev{C_1C_2}$ is the corresponding signed sum of the sixteen correlators.

\section{Results}

For $s=(x,y,x',y')\in\F^4$ put $m_s=\alpha_x\beta_y\alpha'_{x'}\beta'_{y'}$ and
\[
\varphi(s)=x+xx'+xy'+yy' .
\]
Define four functions on $\Omega$,
\[
M_k=c_k\sum_{s\in\F^4}(-1)^{\varphi(s)}\chi_k(s)\,m_s,\qquad k=0,1,2,3,
\]
with $c_1=-1$, $c_0=c_2=c_3=1$, and the characters
\[
\chi_0=1,\qquad \chi_1(s)=(-1)^{y+x'+y'},\qquad \chi_2(s)=(-1)^{x+y'},\qquad \chi_3(s)=(-1)^{x+y+x'} .
\]
Lemma~\ref{l:bent} shows $M_k=\pm4$ at every point of $\Omega$. Let $S_k=\{\omega: M_k(\omega)=-4\}$.
Each $S_k$ has $128$ points, because reversing both outcomes of one party changes the sign of $M_k$.

\begin{theorem}[identity]\label{t:identity}
Let $P$ be a behaviour and $D$ a bilinear form on $W$ satisfying (i) and (ii) of
Definition~\ref{d:jqm}. Then
\[
\sum_{k=0}^{3}D(\ind{S_k},\ind{S_k})\;=\;2-\tfrac14\,\ev{C_1C_2}.
\]
More generally, for any real function $P$ on events and any bilinear $D$ on $W$ with (i) and (ii), the
sum equals $2D(1,1)-\frac14 D(1,C_1C_2)$, and $D(1,1)$ is the total of $P$ over the events of any one
choice of boxes. Pointwise on $\Omega$, $\sum_k\ind{S_k}=2-C_1C_2/4$.
\end{theorem}

\begin{remark}
The four sets are the images of $S_0$ under a group of relabellings. For a party $j$ let $\tau_j$ swap
the two boxes of $j$ and flip the outcome of box $1$ of its partner in the same copy; each $\tau_j$
preserves $\PR\otimes\PR$ and local orthogonality. Then $M_1=M_0\circ\tau_{A_1}$,
$M_2=M_0\circ\tau_{B_1}=M_0\circ\tau_{A_2}$ and $M_3=M_0\circ\tau_{B_2}$, the group generated by the
four $\tau_j$ has order $64$, and $S_0$ has exactly these four images under it. The proof below does not
use the group.
\end{remark}

\begin{corollary}[bound]\label{c:bound}
If a behaviour $P$ has a joint quantum measure on the two-copy scenario, then
$-8\le\ev{C_1C_2}\le 8$. In particular, if $P=P_1\otimes P_2$ with CHSH values $S_1,S_2$ then
$|S_1S_2|\le 8$, and two copies of one box with CHSH value $S$ have a joint quantum measure only if
$|S|\le2\sqrt2$.
\end{corollary}

Positivity is used on eight sets only: the four $S_k$ for the upper bound and their complements for the
lower.

\begin{theorem}[isotropic boxes]\label{t:iso}
$\PR_q\otimes\PR_q$ has a joint quantum measure on the two-copy scenario if and only if $q\le1/\sqrt2$.
\end{theorem}

\begin{theorem}[local orthogonality, exact]\label{t:lo}
$\PR_q\otimes\PR_q$ obeys Consistent Exclusivity on the two-copy scenario if and only if
$q\le q^*:=(\sqrt{10}-1)/3\approx0.7208$. An extremal set of pairwise orthogonal events has $10$
members, of which $5$, $2$ and $3$ are good on both copies, on one, on neither, with total probability
$(10+4q+6q^2)/16$. (An event is \emph{good} on a copy if its two outcomes there satisfy the PR
condition.)
\end{theorem}

The threshold is the one found by Sainz et al.\ \cite{Sainz}, who exhibit a ten-term inequality violated
for $q>q^*$ and report from their computations that this is the optimal value for two copies.
Theorem~\ref{t:lo} proves it: no other inequality of $\mathrm{LO}^2$ is violated earlier.

\begin{corollary}[the converse fails]\label{c:converse}
For $1/\sqrt2<q\le(\sqrt{10}-1)/3$ the behaviour $\PR_q\otimes\PR_q$ obeys Consistent Exclusivity and
has no joint quantum measure. Hence the converse of Henson's theorem \cite{Henson} is false.
\end{corollary}

Both conditions in Corollary~\ref{c:converse} are taken on the same scenario, with the same list of
measurements. Henson's theorem is stated for every scenario; its converse fails on this one.

\section{Proof of Theorem~\ref{t:identity}}

Throughout, $D$ is a bilinear form on $W$ with (i) and (ii), and $P$ is any real function on events.
For a party $j$ let $W_{-j}\subseteq W$ be the span of the indicators of joint outcomes of the other
three parties (one box each).

\begin{lemma}[a shared box drops out]\label{l:rule}
Let $\alpha$ be the outcome of one box of party $j$, and $F,G\in W_{-j}$. Then
$D(\alpha F,\alpha G)=D(F,G)$. Moreover $D(1,\ind e)=D(\ind e,1)=P(e)$ for every event $e$, and
$D(1,1)=\sum_{e\in c}P(e)$ for the set $c$ of the $16$ events of any one choice of boxes.
\end{lemma}

\begin{proof}
Let $\pi=(1+\alpha)/2$ be the indicator of outcome $+1$ of that box. $\pi F$ is a combination of events
in which that box has outcome $+1$, and $(1-\pi)G$ of events in which it has outcome $-1$; every such
pair is locally orthogonal, so $D(\pi F,(1-\pi)G)=0$ by (ii), and likewise $D((1-\pi)F,\pi G)=0$. Since
$\alpha=\pi-(1-\pi)$ and $1=\pi+(1-\pi)$, both $D(\alpha F,\alpha G)$ and $D(F,G)$ equal
$D(\pi F,\pi G)+D((1-\pi)F,(1-\pi)G)$. For the second statement write $1=\sum_{e'\in c}\ind{e'}$ with
$c$ the events on the boxes of $e$; every $e'\ne e$ in $c$ is locally orthogonal to $e$, so
$D(1,\ind e)=D(\ind e,\ind e)=P(e)$ by (ii) and (i), and summing over $e\in c$ gives $D(1,1)$.
\end{proof}

We write $\ev F=D(1,F)$ for $F\in W$. By Lemma~\ref{l:rule} it extends $P$ linearly, $D(F,1)=D(1,F)$,
and for a product of four outcomes it is the correlator of Section~\ref{s:setting}, since such a product
is a signed sum of the sixteen events on its boxes.

\begin{lemma}\label{l:bent}
$M_k=\pm4$ at every point of $\Omega$, for each $k$. Hence $\ind{S_k}=\tfrac12(1-M_k/4)$.
\end{lemma}

\begin{proof}
Fix $\omega\in\Omega$ and write $\alpha_1=\alpha_0(-1)^{t_1}$, $\beta_1=\beta_0(-1)^{t_2}$,
$\alpha'_1=\alpha'_0(-1)^{t_3}$, $\beta'_1=\beta'_0(-1)^{t_4}$. Then $m_s=m_0(-1)^{\ell(s)}$ with
$\ell(s)=t_1x+t_2y+t_3x'+t_4y'$ linear, and $M_k=\pm m_0\sum_s(-1)^{Q(s)}$ where $Q$ is $\varphi$ plus
a linear function. The polar form of $Q$ is that of $\varphi$,
\[
B(s,d)=\varphi(s+d)+\varphi(s)+\varphi(d)=x\,d_{x'}+x'd_x+x\,d_{y'}+y'd_x+y\,d_{y'}+y'd_y ,
\]
and it is non-degenerate: $B(e_x,d)=d_{x'}+d_{y'}$, $B(e_y,d)=d_{y'}$, $B(e_{x'},d)=d_x$,
$B(e_{y'},d)=d_x+d_y$ vanish together only for $d=0$. Therefore
$\bigl(\sum_s(-1)^{Q(s)}\bigr)^2=\sum_d(-1)^{Q(d)}\sum_s(-1)^{B(s,d)}=16$.
\end{proof}

\begin{lemma}[which cross terms survive]\label{l:survive}
For every character $\chi$ of $\F^4$,
\[
\sum_{s,s'}(-1)^{\varphi(s)+\varphi(s')}\chi(s)\chi(s')\,D(m_s,m_{s'})
=16\,D(1,1)+\sum_{d\in\{0111,\,1110,\,1111\}}\chi(d)\,T_d ,
\]
where $d$ is written as $(d_x,d_y,d_{x'},d_{y'})$ and the numbers $T_d$ do not depend on $\chi$.
\end{lemma}

\begin{proof}
Put $d=s+s'$, the set of parties whose box differs between $s$ and $s'$. By Lemma~\ref{l:rule}, applied
once for each party outside $d$, $D(m_s,m_{s'})=D(m_{s|d},m_{s'|d})$, where $m_{s|d}$ is the product of
the outcome variables of the parties in $d$ only. This depends on $d$ and on the restriction of $s$ to
$d$; for $d=0$ it is $D(1,1)$. The sign is $\chi(d)(-1)^{\varphi(s)+\varphi(s+d)}=
\chi(d)(-1)^{\varphi(d)}(-1)^{B(s,d)}$, and $B(s,d)$ is additive in $s$. Summing over the boxes of a
party $j\notin d$ therefore gives zero unless $B(e_j,d)=0$. So a nonzero $d$ contributes only if
$B(e_j,d)=0$ for every $j\notin d$, that is:
$d_{x'}+d_{y'}=0$ if $d_x=0$; $d_{y'}=0$ if $d_y=0$; $d_x=0$ if $d_{x'}=0$; $d_x+d_y=0$ if $d_{y'}=0$.
Going through the fifteen nonzero $d$, exactly $0111$, $1110$ and $1111$ pass. For $d=0$ the sum is over
$16$ values of $s$, each term $D(1,1)$.
\end{proof}

\begin{proof}[Proof of Theorem~\ref{t:identity}]
By Lemma~\ref{l:bent},
\[
D(\ind{S_k},\ind{S_k})=\tfrac14\Bigl[D(1,1)-\tfrac14\bigl(D(1,M_k)+D(M_k,1)\bigr)+\tfrac1{16}D(M_k,M_k)\Bigr].
\]
\emph{Quadratic terms.} By Lemma~\ref{l:survive} with $\chi=\chi_k$ (and $c_k^2=1$),
$\sum_kD(M_k,M_k)=64\,D(1,1)+\sum_dT_d\sum_k\chi_k(d)$, and
\[
\sum_k\chi_k(0111)=1-1-1+1,\qquad\sum_k\chi_k(1110)=1+1-1-1,\qquad\sum_k\chi_k(1111)=1-1+1-1 ,
\]
all zero. So $\sum_kD(M_k,M_k)=64\,D(1,1)$, whatever the values $T_d$.

\emph{Linear terms.} With $a=x+y'$ and $b=x+y+x'$ one has $y+x'+y'=a+b$ modulo $2$, so
\[
\sum_kc_k\chi_k(s)=1-(-1)^{a+b}+(-1)^a+(-1)^b=2(-1)^{ab},
\]
and $\varphi(s)+ab=xy+x'y'$ modulo $2$. Hence $\sum_kM_k=2\sum_s(-1)^{xy+x'y'}m_s=2\,C_1C_2$, which is the
pointwise identity, and $\sum_kD(1,M_k)=\sum_kD(M_k,1)=2\,D(1,C_1C_2)$.

Adding up: $\sum_kD(\ind{S_k},\ind{S_k})=\tfrac14\bigl[4D(1,1)-D(1,C_1C_2)+4D(1,1)\bigr]
=2D(1,1)-\tfrac14D(1,C_1C_2)$. For a behaviour, $D(1,1)=1$ by Lemma~\ref{l:rule}.
\end{proof}

\section{Proofs of the other results}

\begin{proof}[Proof of Corollary~\ref{c:bound}]
By (iii) each $D(\ind{S_k},\ind{S_k})\ge0$, so $\ev{C_1C_2}\le8$. The complements have indicators
$\tfrac12(1+M_k/4)$, and the same computation gives
$\sum_kD(1-\ind{S_k},1-\ind{S_k})=2+\tfrac14\ev{C_1C_2}\ge0$. For a product behaviour
$\ev{C_1C_2}=S_1S_2$.
\end{proof}

\begin{proof}[Proof of Theorem~\ref{t:iso}]
\emph{Only if.} $\ev{C_1C_2}=16q^2$, which exceeds $8$ when $q>1/\sqrt2$.

\emph{If.} Let $q\le1/\sqrt2$. $\PR_q$ is obtained from projective $\pm1$ measurements on two qubits
(those that give Tsirelson's bound on the maximally entangled state) in the state
$\rho_1=t\,|\Phi\rangle\langle\Phi|+(1-t)\,\mathbf 1/4$ with $t=\sqrt2\,q$. Take two copies, state
$\rho=\rho_1\otimes\rho_1$, and for party $j$ let $\Pi^j_{b,a}$ be the projector for outcome $a$ of box
$b$, acting on that party's qubit. For $\gamma\in\Omega$ put
\[
V_\gamma=\bigotimes_j\Pi^j_{0,\gamma_{j0}}\Pi^j_{1,\gamma_{j1}},\qquad
D(\ind\gamma,\ind{\gamma'})=\mathrm{Re}\,\mathrm{tr}\bigl(\rho\,V_\gamma^\dagger V_{\gamma'}\bigr),
\]
extended bilinearly. $D$ is a real symmetric positive semidefinite form on all functions on $\Omega$,
so (iii) holds. For an event $e$, summing $V_\gamma$ over $\gamma\in e$ removes, for each party, the
projector of the box that $e$ does not open, because the projectors of one box sum to the identity:
$V_e:=\sum_{\gamma\in e}V_\gamma=\bigotimes_j\Pi^j_{b_j,a_j}$. Hence $D(\ind e,\ind e)=\mathrm{tr}(\rho V_e)
=P(e)$, which is (i), and if $e,e'$ are locally orthogonal then some party contributes two orthogonal
projectors of one box to $V_e^\dagger V_{e'}$, which is therefore zero: (ii). The fine-grained outcomes
of any branching measurement are pairwise locally orthogonal events, so for its coarse-grained outcomes
$X,Y$ one gets $D(\ind X,\ind Y)=P(X\cap Y)$: $D$ is a strongly positive joint decoherence functional
in the sense of \cite{DHW} for the scenario with branching.
\end{proof}

\begin{proof}[Proof of Theorem~\ref{t:lo}]
An event that is good on both copies, on one, on neither has probability $(1+q)^2/16$, $(1-q^2)/16$,
$(1-q)^2/16$. A set of pairwise orthogonal events with $a$, $b$, $c$ members of the three kinds has
total $[a(1+q)^2+b(1-q^2)+c(1-q)^2]/16$.

\emph{Only if.} There are sets of type $(a,b,c)=(5,2,3)$ ($18{,}944$ of them); their total
$(10+4q+6q^2)/16$ exceeds $1$ exactly when $3q^2+2q-3>0$, that is $q>q^*$.

\emph{If, at $q^*$.} Orthogonal events are disjoint $16$-point subsets of $\Omega$, so a set of pairwise
orthogonal events has at most $16$ members; at most $5$ of them are good on both copies \cite{Sainz}.
With $k\le3$ such members the total is at most $[k(1+q)^2+(16-k)(1-q^2)]/16<1$ at $q=q^*$. For $k=4$
and $k=5$ the maximum is found by exhaustive search over the extensions of the $18{,}160$ and $640$
sets of $k$ pairwise orthogonal good-good events. At $q^*$ the three probabilities, times $144$, are
$14+4\sqrt{10}$, $-2+2\sqrt{10}$, $26-8\sqrt{10}$, so the excess of a set's total over $1$, times $144$,
is $p+r\sqrt{10}$ with integers $|p|\le560$, $|r|\le128$, and a nonzero number of this form has modulus
at least $1/(560+128\sqrt{10})>10^{-3}$. The search is carried out in integers with the weights scaled by
$10^9$ and rounded to the nearest integer, so that a set's rounding error is at most $16\cdot\frac12=8$.
It finds no set whose scaled total exceeds $144\cdot10^9+8$, whereas a set with positive excess would
reach $144\cdot10^9+10^6-8$. Types
$(5,2,3)$ and $(4,8,4)$ have excess exactly $0$.

\emph{If, below $q^*$.} Let $t=q/q^*\in[0,1]$. Then $\PR_q=t\,\PR_{q^*}+(1-t)U$, so $\PR_q\otimes\PR_q$
is a convex combination of $\PR_{q^*}\otimes\PR_{q^*}$, $\PR_{q^*}\otimes U$, $U\otimes\PR_{q^*}$ and
$U\otimes U$, and CE is a family of linear inequalities. $U$ is a mixture of deterministic boxes. For
$B\otimes\delta$ with $\delta$ deterministic on copy 2, the events of nonzero probability carry the
copy-2 outcomes that $\delta$ dictates, so two of them cannot be orthogonal through a party of copy 2; a
set of pairwise orthogonal such events restricts to distinct pairwise orthogonal events of copy 1, whose
total under the no-signalling box $B$ is at most $1$, since in a bipartite scenario local orthogonality
is implied by no-signalling \cite{Sainz}. Likewise with the copies exchanged.
\end{proof}

Corollary~\ref{c:converse} is Theorems~\ref{t:iso} and \ref{t:lo} together.

\section{What the result needs}\label{s:needs}

Call a wing one experimenter's two parties: $\{A_1,A_2\}$ and $\{B_1,B_2\}$. A symmetric form on $W$
has $81\cdot82/2=3{,}321$ free entries ($\dim W=81$). Table~\ref{tab:rank} gives, for four lists of
measurements, the number of pairs of events that decohere, the rank of conditions (i) and (ii) on those
entries, and whether the four-set sum of Theorem~\ref{t:identity} is fixed by them. (Computed in exact
integer arithmetic; the rank includes the $256$ diagonal conditions, and the sum is fixed exactly when
adding its functional to the conditions does not raise the rank.)

\begin{table}[h]
\centering
\begin{tabular}{lrrl}
measurements & decohering pairs & rank & four-set sum \\ \hline
all branching & $22{,}400$ & $3{,}049$ & fixed \\
branching within each wing only & $8{,}064$ & $2{,}153$ & not fixed \\
branching within each copy only & $8{,}064$ & $2{,}153$ & not fixed \\
no branching & $1{,}920$ & $1{,}241$ & not fixed
\end{tabular}
\caption{The identity needs branching across wings.}\label{tab:rank}
\end{table}

The benchmark needs the same premise. Every one of the $18{,}944$ extremal sets of
Theorem~\ref{t:lo} contains at least $18$ pairs (of its $45$) that are not outcomes of a common product
of one branching measurement per wing. With such product measurements only, at most four good-good
events are pairwise exclusive, and Consistent Exclusivity on two copies excludes no isotropic box, not
even $\PR\otimes\PR$. This is the two-copy face of a known fact: in a bipartite scenario local
orthogonality is equivalent to no-signalling, and the two-copy violation appears when the copies are
``shared among $2k$ parties'' \cite{Sainz}. The class of pairs exclusive through one branching
measurement per wing is not the bipartite one: $3{,}072$ of its pairs are not orthogonal when each wing
is treated as one party with four boxes, and $6{,}144$ pairs are orthogonal in that sense without
belonging to it.

We have no sharp bound on $\langle C_1C_2\rangle$ for joint quantum measures when branching is allowed
within wings only. The result gives Tsirelson's bound for CHSH on two copies; it does not show that JQM
together with closure under composition implies strong positivity, which remains open \cite{DHW}. Three
or more copies, other Bell inequalities and other scenarios are open.

\section{Earlier work}\label{s:earlier}

Henson \cite{Henson} proved that JQM implies Consistent Exclusivity and excluded two copies with CHSH
value above $2.883$. Dowker, Henson and Wallden \cite{DHW} introduced JQM and its strongly positive
form, noted that JQM is not closed under composition and that $\PR\otimes\PR$ is not in JQM, and asked
whether strong positivity follows from JQM with closure under composition. Bo\"es and Navascu\'es
\cite{BN} showed that, under a product rule for the decoherence functionals of independent systems,
closure under composition for all numbers of copies forces strong positivity; we assume no product rule,
use two copies only, and exclude every joint quantum measure of the two-copy behaviour. Niestegge
\cite{Niestegge} derived Tsirelson's bound from the absence of third-order interference in a
quantum-logical framework with further axioms; whether that definition of third-order interference is
equivalent to Sorkin's is left open in \cite{Henson}.

Sainz et al.\ \cite{Sainz} obtained the value $(\sqrt{10}-1)/3$ and write of their five-event set that
it ``is the only violation of $\mathrm{LO}^2$ by the PR-box up to symmetry''. We find that the $640$
sets of five pairwise orthogonal good-good events fall into two classes, of $512$ and $128$ sets. In
each set exactly two of the ten pairs are orthogonal through two parties; those two pairs share an
event in the $512$ sets and do not in the $128$. No relabelling of parties, boxes and outcomes maps one
class to the other, and no automorphism of the orthogonality graph does: for a set in the $128$-class
there are events outside it orthogonal to exactly one of its five members, and for a set in the
$512$-class there are none. Both sets printed in \cite{Sainz} are in the $512$-class. Each set of that
class extends to a set of type $(5,2,3)$ in $30$ ways, each set of the other in $28$.

\section{Machine-checked statements}

The following are formalised in Lean~4 with Mathlib, with no axioms beyond the standard three and with
finite facts evaluated in the kernel. In the formalisation the four sets are defined as the images of
$S_0$ under the relabellings $\tau_j$ of the remark after Theorem~\ref{t:identity}; that these are the
sets $S_k$ defined by the characters $\chi_k$ is the computation in that remark and is not formalised.
Formalised: $M_0=\pm4$ and $|S_0|=128$; that $S_0$ has exactly four images under the relabelling group
and that the group has order $64$; the pointwise identity; the first statement of Lemma~\ref{l:rule} and
the value of $D(1,1)$; Theorem~\ref{t:identity} in its general form (bilinear $D$ on $W$, not assumed
symmetric, any real $P$); both bounds of Corollary~\ref{c:bound} for a behaviour; the ``only if'' half
of Theorem~\ref{t:iso}; and that each event is locally orthogonal to $175$ others. Not
formalised: the product sentences of Corollary~\ref{c:bound}, the ``if'' half of Theorem~\ref{t:iso},
Theorem~\ref{t:lo}, and the counts and ranks of Sections~\ref{s:needs} and \ref{s:earlier}, which are computed by the
scripts filed with this paper.

\begin{thebibliography}{9}
\bibitem{Sorkin} R.~D. Sorkin, Quantum mechanics as quantum measure theory, \emph{Mod.\ Phys.\ Lett.}\
A9 (1994) 3119--3128, arXiv:gr-qc/9401003.
\bibitem{DHW} F.~Dowker, J.~Henson and P.~Wallden, A histories perspective on characterizing quantum
non-locality, \emph{New J.\ Phys.}\ 16 (2014) 033033, arXiv:1311.6287.
\bibitem{Henson} J.~Henson, Bounding quantum contextuality with lack of third-order interference,
arXiv:1406.3281.
\bibitem{Sainz} A.~B. Sainz, T.~Fritz, R.~Augusiak, J.~Bohr Brask, R.~Chaves, A.~Leverrier and
A.~Ac\'in, Exploring the local orthogonality principle, arXiv:1311.6699.
\bibitem{BN} P.~Bo\"es and M.~Navascu\'es, Composing decoherence functionals, arXiv:1609.09723.
\bibitem{Niestegge} G.~Niestegge, Three-slit experiments and quantum nonlocality, \emph{Found.\ Phys.}\
43 (2013) 805--812, arXiv:1104.0091.
\bibitem{PR} S.~Popescu and D.~Rohrlich, Nonlocality as an axiom, \emph{Found.\ Phys.}\ 24 (1994) 379.
\bibitem{Tsirelson} B.~S. Cirel'son, Quantum generalizations of Bell's inequality, \emph{Lett.\ Math.\
Phys.}\ 4 (1980) 93--100.
\end{thebibliography}

\end{document}
